\documentclass[10pt]{article}
\usepackage[T1]{fontenc}
\usepackage[utf8]{inputenc}
\usepackage[french]{babel}
\usepackage{amsmath,amssymb,mathrsfs,mathpazo}
\usepackage{geometry}
\geometry{a4paper,margin=2cm}
\usepackage{graphicx,array,tabularx,tabularray,colortbl,multirow,multicol,booktabs,enumitem,pifont,ulem,fancybox,diagbox,cancel}
\usepackage{pstricks,pstricks-add,pst-plot,pst-tree,pst-node,pst-text,pst-func,pst-math,pst-bspline,pst-3dplot,pst-eucl,pst-solides3d}
\let\pscalculate\undefined
\usepackage{pst-bezier}
\usepackage{tikz,pgfplots,tkz-tab}
\usetikzlibrary{arrows,arrows.meta,patterns,calc,positioning,trees,decorations.markings,shapes.geometric,shapes.misc,shapes.arrows,intersections,decorations.pathmorphing,decorations.pathreplacing,backgrounds,shadows,babel}
\pgfplotsset{compat=1.18}
\usepgfplotslibrary{fillbetween}
\usepackage[np]{numprint}
\usepackage{listings,makecell,pas-tableur,twoopt}
\let\barre\undefined
\usepackage[pstricks]{tablvar}
\definecolor{darkgreen}{RGB}{0,100,0}
\providecommand{\red}{\color{red}}
\providecommand{\blue}{\color{blue}}
\providecommand{\green}{\color{green}}
\providecommand{\black}{\color{black}}
\providecommand{\white}{\color{white}}
\providecommand{\gray}{\color{gray}}
\providecommand{\cyan}{\color{cyan}}
\providecommand{\magenta}{\color{magenta}}
\providecommand{\yellow}{\color{yellow}}
\providecommand{\violet}{\color{violet}}
\providecommand{\orange}{\color{orange}}
\providecommand{\euro}{\texteuro}
\providecommand{\EUR}{\texteuro}
\providecommand{\ieme}{\textsuperscript{e}}
\providecommand{\degres}{\ensuremath{^\circ}}
\providecommand{\degree}{\ensuremath{^\circ}}
\providecommand{\vv}[1]{\overrightarrow{#1}}
\providecommand{\Coord}[1]{#1}
\providecommand{\decosix}{}
\providecommand{\starredbullet}{\ensuremath{\bullet}}
\pagestyle{empty}
\setlength{\parindent}{0pt}
\usepackage{adjustbox}\DeclareUnicodeCharacter{2236}{:}\DeclareUnicodeCharacter{2190}{\ensuremath{\leftarrow}}
\providecommand{\R}{\mathbb{R}}
\providecommand{\N}{\mathbb{N}}
\providecommand{\Z}{\mathbb{Z}}
\providecommand{\Q}{\mathbb{Q}}
\providecommand{\D}{\mathbb{D}}
\providecommand{\C}{\mathbb{C}}
\providecommand{\e}{{\mathrm{e}}}
\providecommand{\vect}[1]{\overrightarrow{#1}}
\providecommand{\vectt}[1]{\overrightarrow{\mathrm{#1}}}
\providecommand{\vvt}[1]{\overrightarrow{\mathrm{#1}}}
\providecommand{\np}[2][]{{#2~#1}}
\providecommand{\vv}[1]{{\overrightarrow{#1}}}
\renewcommand{\d}{{\mathrm{d}}}
\providecommand{\degres}{{^\circ}}
\providecommand{\degre}{{^\circ}}
\providecommand{\blue}{\color{blue}}
\providecommand{\red}{\color{red}}
\providecommand{\enskip}{\quad}
\providecommand{\llbracket}{\lbrack\!\lbrack}
\providecommand{\rrbracket}{\rbrack\!\rbrack}
\providecommand{\euro}{{\text{€}}}
\providecommand{\no}{{\text{n°}}}
\providecommand{\interval}[2]{{[#1;#2]}}
\providecommand{\intervFO}[2]{{[#1;#2[}}
\providecommand{\dotfill}{\ldots}

\providecommand{\ds}{\displaystyle}
\providecommand{\textding}[1]{#1}
\providecommand{\Oij}{\(\mathrm{O};\vec{\imath},\vec{\jmath}\)}
\providecommand{\Oijk}{\(\mathrm{O};\vec{\imath},\vec{\jmath},\vec{k}\)}
\providecommand{\Ouv}{\(\mathrm{O};\vec{u},\vec{v}\)}
\providecommand{\intervFF}[2]{\ensuremath{[#1;#2]}}
\renewcommand{\euro}{\ifmmode\text{\texteuro}\else\texteuro\fi}
\DeclareUnicodeCharacter{00B0}{\ifmmode^{\circ}\else\textdegree\fi}
\DeclareUnicodeCharacter{20AC}{\ifmmode\text{\texteuro}\else\texteuro\fi}
% Énoncés : Éducation nationale. Transcription et corrigés : APMEP, Denis Vergès et contributeurs. Découpage et classement : CoopMaths, Sébastien Lozano, Éric Elter et contributeurs. Adaptation de la présentation : Maths Explorer.
% CC BY-SA : https://coopmaths.fr/www/about/
% https://coopmaths.fr/alea/static/bac/2022/tex/bac_2022_05_sujet2_etrangers_1_cor.tex
\begin{document}
\noindent\textit{APMEP / CoopMaths. CC BY-SA. Présentation adaptée par Maths Explorer.}\par\medskip
\medskip

\begin{enumerate}
	\item Soit $f$ la fonction continue et dérivable sur $\R$, définie par $x \longmapsto \dfrac{x}{\text{e}^x}$.
	
	En utilisant la formule de la dérivée d'un quotient, on peut écrire :
	
	$\forall x \in \R$, $f'(x)=\dfrac{1\times \e^{x}-x\times \e^x}{\left(\e^x\right)^2}=\dfrac{(1-x)\e^x}{\e^{2x}}=\dfrac{1-x}{\e^x}=(1-x)\e^{-x}$. \hfill \fbox{\textbf{Réponse c}}
	\item La courbe représentative donnée est celle de $f''$. Avec la précision permise par ce graphique, on peut affirmer que $f''$ est positive sur $[-3~;~-1]$, négative sur $[-1~;~1]$ et s'annule pour $x= - 1$. 
	
Cela signifie donc que $f'$ est croissante sur $[-3~;~-1]$ puis décroissante sur $[-1~;~1]$. Donc la fonction $f'$ admet un maximum en $x = -1$. \hfill \fbox{\textbf{Réponse d}}
	\item Les fonctions $F$ proposées sont continues et dérivables sur $\R$. Dérivons chacune d'entre-elles. $\forall x \in \R$, 
	\begin{enumerate}
		\item $\begin{gathered}F'(x)=-\dfrac{1}{6}\left(3x^2\e^{-x^2}+(x^3+1)\times -2x\e^{-x^2}\right)=-\dfrac{1}{6}\e^{-x^2}(3x^2-2x^3-2x)\\=-\dfrac{1}{6}\e^{-x^2}(-2x^3+3x^2-2x)\end{gathered}$
		\item $F'(x)=-\dfrac{1}{4}\left(4x^3 \times \e^{-x^2}+x^4\times -2x\e^{-x^2}\right)=-\dfrac{1}{4}\e^{-x^2}(4x^3-2x^5)$
		\item $F'(x)=-\dfrac{1}{2}\left(2x\times \e^{-x^2}+(x^2+1)\times -2x\e^{-x^2}\right)=-\dfrac{1}{2}\e^{-x^2}(2x-2x^3-2x)=x^3\e^{-x^2}$
		\item $F'(x)=(6x-8x^3)\times \e^{-x^2}+(3x^2-2x^4)\times -2x\e^{-x^2}= (4x^5-14x^3+6x)\e^{-x^2} $
	\end{enumerate}
\hfill \fbox{\textbf{Réponse c}}
	\item $\forall x \in \R$, $f(x)=\dfrac{\e^x\left(1+\dfrac{1}{\e^x}\right)}{\e^x\left(1-\dfrac{1}{\e^x}\right)}=\dfrac{1+\e^{-x}}{1-\e^{-x}}$. Or $\displaystyle\lim\limits_{x \rightarrow +\infty} \e^{-x}=0$ donc $\displaystyle\lim\limits_{x \rightarrow +\infty} f(x)=1$. \hfill \fbox{\textbf{Réponse b}}
	\item La fonction $f$ est continue sur $\R$, donc elle admet des primitives notées $F$. De plus $f(x)$ est de la forme $u'(x)\e^{u(x)}$ avec $u(x)=2x+1$. En remarquant que $f(x)=\dfrac{1}{2} \times 2\e^{2x+1}$, on peut alors écrire que pour tout réel $x$, $F(x)=\dfrac{1}{2}\e^{2x+1}+k$ avec $k\in \R$. Sachant que $F(0)=1$ alors
	
	 $\dfrac{1}{2}\e^1+k=1 \iff k=1-\dfrac{1}{2}\e$. Donc $\forall x \in \R$, $F(x)=\dfrac{1}{2}\e^{2x+1}+1-\dfrac{1}{2}\e$. \hfill \fbox{\textbf{Réponse c}}
	 \item Avec la précision permise par le graphique, on peut affirmer que la fonction $f$ est concave sur $[-2~;~1]$, convexe sur $[1~;~4]$ et $f''(1)=0$ (point d'inflexion). $f''$est donc négative sur $[-2~;~1]$, positive sur $[1~;~4]$ s'annulant en 1. \hfill \fbox{\textbf{Réponse a}}
\end{enumerate}

\pagebreak

\end{document}
